\documentclass[a4paper]{article}
\usepackage{amsmath,amssymb}
\usepackage{geometry}
\geometry{left=25mm,right=25mm,top=25mm,bottom=25mm}
\usepackage{hyperref}

\title{De Broglie Waves as Dual-Rotor Phase Lag:\\ Quantum Mechanics as Geometry of Particle-Intrinsic Dual Spacetime}

\author{https://github.com/hypernumbernet}
\date{\today}

\begin{document}
\maketitle

\begin{abstract}
Every massive particle in double spacetime theory carries a dual rotor, already an element of \(\mathrm{SU}(2)\) on \(\mathbb{C}^2\). Three invariants of that rotor must be kept apart. The character \(\operatorname{Tr} R(\beta)=2\cos(\lVert\beta\rVert/2)\) is real because it is the sum of opposite half-angle phases on the two computational spinors. The complex matrix element \(\langle\uparrow\rvert R(\beta)\lvert\uparrow\rangle=\cos(\theta/2)-i n_z\sin(\theta/2)\) is a pure phase along the observer axis and a real cosine in the equatorial plane; its modulus is the polar tilt of the dual axis. A scalar \(\mathrm{U}(1)\) factor identified with the de Broglie phase is a third, distinct quantity. Observer and particle are compared by the Hilbert pairing of their dual-rotor images, equal to the matrix element of a relative dual rotor that itself remains in \(\mathrm{SU}(2)\). That relative rotor is a product in the group; its character records the alignment of the two dual axes. That same matrix element is the overlap of rest-frame Dirac solutions that differ by an internal dual rotation. The written opposite-sign particle action makes the mismatch free, so lag may accumulate linearly; a harmonic model exists only for same-sign kinetics, at frequency \(\sqrt{2m}\), which is not \(E/\hbar\). The Dirac equation in momentum space is the first-order factor of the Minkowski quadratic in \(\mathrm{Cl}(3,1)\). Schr\"odinger and Klein--Gordon dynamics, collapse, and the unification of gravity with quantum non-locality remain geometric readings. The phase \(J\) of this note is an oscillator diagnostic, not the parameter-space Killing quadratic of the main paper.
\end{abstract}

\section{Introduction}

De Broglie's 1924 hypothesis that every massive particle is accompanied by a wave of wavelength \(\lambda=h/p\) remains one of the deepest empirical facts in physics~\cite{debroglie1924}. Standard quantum mechanics places the wave function in an abstract Hilbert space. General relativity places gravity in the curvature of a shared continuum. Neither construction says, from the algebra carried by a single massive particle, \emph{what} the matter wave is.

Double spacetime theory supplies a candidate~\cite{dst-main}. Every massive particle carries two compactified Minkowski frames inside the 16-real-dimensional biquaternion algebra \(\mathbb{H}\otimes\mathbb{H}\cong\mathrm{Cl}(3,1)\). Gravity and inertia are torsional mismatch between the usual rotor \(R_{\mathrm{usual}}\) and the dual rotor \(R_{\mathrm{dual}}\). The dual rotor, represented on \(\mathbb{C}^2\) by
\[
R(\beta)=\exp\Bigl(\sum_{a=1}^3\tfrac{\beta_a}{2}(-i\sigma_a)\Bigr),
\]
is unitary of determinant \(1\), hence an element of \(\mathrm{SU}(2)\). Compactness of the dual sector is the \(4\pi\) period of that group.

The question is which invariant of \(R(\beta)\) is the de Broglie phase. The character, the spinor matrix element, and a scalar \(\mathrm{U}(1)\) factor built from the mismatch angle are three different quantities. Conflating them produces a complex wave function out of a real trace, or a Compton frequency out of a Lagrangian that does not oscillate. The geometry is clearer when the three are named separately.

\section{Dual rotors}

The complete rotor of a massive particle is written
\[
R_{\mathrm{total}}
  =\exp\Bigl[\sum_{a=1}^3\bigl(\tfrac{\theta_a}{2}i\Gamma_a+\tfrac{\phi_a}{2}\Gamma_a\bigr)\Bigr],
\qquad \Gamma_1=I,\;\Gamma_2=J,\;\Gamma_3=K,
\]
where \(i\Gamma_a\) generate boosts (square to \(+1\)) and \(\Gamma_a\) generate dual-sector rotations (square to \(-1\)). Same-axis usual and dual generators commute, so a one-axis exponential factorises. Distinct-axis generators need not commute; an unrestricted six-parameter factorisation \(R_{\mathrm{usual}}R_{\mathrm{dual}}\) is not justified by those commutators.

The relative rotor \(\Omega=R_{\mathrm{usual}}^{\dagger}R_{\mathrm{dual}}\) equals the identity if and only if the two rotors coincide. Vanishing of the parameter-space scalar \(J\) does not force \(\Omega=1\): a balanced massive seed has \(J=0\) with a nontrivial relative rotor.

On \(\mathbb{C}^2\) the dual rotor is the Rodrigues matrix
\[
R(\beta)=\cos(\lVert\beta\rVert/2)\,I+\sin(\lVert\beta\rVert/2)\,\widehat{\Gamma}(\beta),
\]
with \(\widehat{\Gamma}(\beta)^2=-1\) whenever \(\beta\neq 0\). Usual-sector rotors are the hyperbolic twin: Hermitian of determinant \(1\), and not unitary for a nonzero boost.

\section{Three invariants}

The character of the dual rotor is the normalised trace
\[
\frac12\operatorname{Tr} R(\beta)=\cos(\lVert\beta\rVert/2).
\]
It is real, and it is real for a geometric reason. Let \(\lvert\uparrow\rangle\) and \(\lvert\downarrow\rangle\) be the computational eigenvectors of \(\sigma_z\) in the observer's dual frame. The trace is their sum,
\[
\operatorname{Tr} R(\beta)
  =\langle\uparrow\rvert R(\beta)\lvert\uparrow\rangle
  +\langle\downarrow\rvert R(\beta)\lvert\downarrow\rangle,
\]
so the half-trace is the real part of the spin-up matrix element. Along the observer axis the two amplitudes are opposite half-angle phases,
\[
\langle\uparrow\rvert R(\theta e_z)\lvert\uparrow\rangle=e^{-i\theta/2},\qquad
\langle\downarrow\rvert R(\theta e_z)\lvert\downarrow\rangle=e^{+i\theta/2},
\]
and their sum is \(2\cos(\theta/2)\). A pairing built from this trace cannot be a complex scalar wave function: the two \(\mathrm{U}(1)\) phases have already been added, and the imaginary part has cancelled. What remains is the \(\mathrm{SU}(2)\) character of the dual rapidity, the half-angle cosine of a rotation on the three-sphere of unit dual rotors.

The finer invariant is a single spinor matrix element. With \(\theta=\lVert\beta\rVert\) and \(\hat n=\beta/\theta\),
\[
\langle\uparrow\rvert R(\beta)\lvert\uparrow\rangle
  =\cos(\theta/2)-i n_z\sin(\theta/2).
\]
Its squared modulus is the polar tilt of the dual axis against the observer,
\[
\bigl\lvert\langle\uparrow\rvert R(\beta)\lvert\uparrow\rangle\bigr\rvert^2
  =\cos^2(\theta/2)+n_z^2\sin^2(\theta/2)
  =1-(n_x^2+n_y^2)\sin^2(\theta/2).
\]
A pure phase, of unit modulus, appears only when that axis is aligned with the observer, or when the rotation is trivial. Two axes display the same geometry cleanly. A dual rotation about the observer axis is the phase \(e^{-i\theta/2}\) already written. A dual rotation in the equatorial plane is a real cosine,
\[
\langle\uparrow\rvert R(\theta e_x)\lvert\uparrow\rangle=\cos(\theta/2).
\]
Amplitude and phase of the matrix element are therefore not independent decorations of a single angle: they record how the dual axis sits relative to the observer. The polar form \(\sqrt{\rho}\,e^{iS/\hbar}\) of textbook wave mechanics is not this \(\mathrm{SU}(2)\) object, and cannot be recovered from \(\arg(\hat n_{\mathcal{O}}\cdot\hat n_A)\), which for two real unit vectors is \(0\) or \(\pi\).

The third quantity is a scalar mismatch angle \(\delta\) between usual and dual rapidities, promoted by hand to a laboratory \(\mathrm{U}(1)\) factor \(e^{i\delta}\). That identification may be read as the de Broglie phase \((Et-\mathbf{p}\cdot\mathbf{x})/\hbar\). It is not the character, and it is not automatically the spinor matrix element. The double cover of \(\mathrm{SU}(2)\) makes the distinction a matrix identity: a \(2\pi\) dual rotation about the observer axis is \(-I\), so the spinor goes to its negative and the half-trace to \(\cos\pi=-1\); a \(4\pi\) rotation is the identity. A conventional scalar matter wave would have completed a single cycle already at \(2\pi\).

\section{Observer overlap}

Let \(\mathcal{O}\) be a macroscopic observer and \(A\) a massive particle. The comparison that lives in the dual Hilbert space is
\[
\langle R(\beta_{\mathcal{O}})\chi\mid R(\beta_A)\chi\rangle
  =\langle\chi\mid R(\beta_{\mathcal{O}})^{\dagger}R(\beta_A)\chi\rangle,
\]
for a reference spinor \(\chi\). Unitarity of each dual rotor makes the left-hand side an inner product of two unitarily transported copies of \(\chi\), and makes the right-hand side the matrix element of the relative dual rotor \(R(\beta_{\mathcal{O}})^{\dagger}R(\beta_A)\). That relative rotor is itself unitary of determinant \(1\), so the pairing never leaves \(\mathrm{SU}(2)\). The modulus is bounded by \(\lVert\chi\rVert^2\).

This pairing is not the parameter-space Killing form. The Killing pairing has signature \((3,3)\). Self-overlap equals a multiple of \(J\) and can be negative; it is not a Born probability. The four labels of dual-spacetime logic are not four orthogonal rays in \(\mathbb{C}^2\).

The inverse of a dual rotor is the dual rotor of the opposite rapidity, \(R(\beta)^{\dagger}=R(-\beta)\). The relative rotor is therefore itself a product of dual rotors, \(R(-\beta_{\mathcal{O}})R(\beta_A)\). Along a common ray the generators commute, so the product collapses to a single dual rotation by the difference of the scales: \(R(s\hat n)^{\dagger}R(t\hat n)=R((t-s)\hat n)\). That is the addition of angles already visible on a shared axis, now read as a comparison rather than a transport.

Off that axis the rapidities do not add. The quaternion table \(\Gamma_a\Gamma_b=-\delta_{ab}+\varepsilon_{abc}\Gamma_c\) expands the Rodrigues product into a scalar and a traceless remainder. The scalar is the character of the composite,
\[
\tfrac12\operatorname{Tr}\bigl(R(\beta)R(\gamma)\bigr)
  =\cos(\theta/2)\cos(\phi/2)-(\hat n\cdot\hat m)\sin(\theta/2)\sin(\phi/2),
\]
with \(\theta=\lVert\beta\rVert\), \(\phi=\lVert\gamma\rVert\), and \(\hat n\), \(\hat m\) the two dual axes. The character of a product is therefore not the product of the characters: it is shifted by the inner product of the axes. For the relative rotor the inverse flips that sign,
\[
\tfrac12\operatorname{Tr}\bigl(R(\beta_{\mathcal{O}})^{\dagger}R(\beta_A)\bigr)
  =\cos(\theta_{\mathcal{O}}/2)\cos(\theta_A/2)+(\hat n_{\mathcal{O}}\cdot\hat n_A)\sin(\theta_{\mathcal{O}}/2)\sin(\theta_A/2).
\]
Aligned axes recover the half-angle cosine of the difference. Orthogonal axes leave only the product of cosines. The pairing of a single computational spinor is still a matrix element, not this character; the character has already added the two opposite phases.

The remainder of the product is traceless. It is a real linear combination of \(\Gamma(\hat n)\), \(\Gamma(\hat m)\), and \(\Gamma(\hat n\times\hat m)\). The cross product is the noncommutative residue: it tilts the axis of the composite and is invisible to the half-trace. Superposition of distinct dual axes is therefore composition in \(\mathrm{SU}(2)\), not addition of complex scalars. A sharp instance is a pair of \(\pi\) dual rotations about perpendicular axes. They compose to a \(\pi\) rotation about the third, opposite order giving the opposite sign, and their group commutator is \(-I\)---the same central inversion as a \(2\pi\) dual rotation about the observer. Non-abelian interference, in this geometry, is not a small correction to a scalar fringe. It is the double cover of \(\mathrm{SU}(2)\).

\section{Mismatch dynamics}

A particle-intrinsic action for the six rapidities is a model, not a deduction from unitarity. Three neighbouring systems have been compared~\cite{dst-dynamic-equation,dst-main}.

The written opposite-sign density
\[
L=\tfrac12\dot\phi^2-\tfrac12\dot\theta^2-\tfrac m2(\phi-\theta)^2
\]
has Euler--Lagrange equations \(\ddot\phi=-m(\phi-\theta)\) and \(\ddot\theta=-m(\phi-\theta)\). The mismatch \(\delta=\phi-\theta\) is free: \(\ddot\delta=0\). Affine lag \(\delta(t)=\delta_0+\nu t\) is allowed. The coefficient \(\nu\) is an initial condition; the mass \(m\) accelerates both angles together and does not fix a de Broglie frequency.

A neighbouring system sometimes written alongside it,
\[
\ddot\phi=m(\phi-\theta),\qquad -\ddot\theta=m(\phi-\theta),
\]
makes the mismatch runaway, \(\ddot\delta=2m\delta\). It is not an oscillator.

An oscillator appears only if both kinetic signs are taken positive. The density \(L=\tfrac12\dot\phi^2+\tfrac12\dot\theta^2-\tfrac m2(\phi-\theta)^2\) yields \(\ddot\delta+2m\delta=0\), with frequency \(\sqrt{2m}\) in units \(\hbar=c=1\). That frequency is not the Compton frequency \(m\), and the model is not the written action.

Linear accumulation of a free mismatch is therefore the only dynamical bridge to a secular phase that the written action actually supplies. Identifying the rate with \(E/\hbar\) remains a reading. Collective averaging of a particle-intrinsic \(J_{\mathrm{osc}}\) over many particles is likewise a reading; it is not an identification with the Weitzenb\"ock density \(T\), and \(\lvert J_{\mathrm{norm}}\rvert\le 1\) is not transplanted here.

The Compton wavelength continues to be readable as a compactification radius of the dual sector, \(\sim\hbar/mc\). The de Broglie wavelength is the laboratory distance over which an identified scalar phase advances by \(2\pi\). Neither length is derived from the trace character.

\section{Dirac equation as first-order form}

The underlying algebra is \(\mathrm{Cl}(3,1)\). A first-order spinor equation is therefore not an extra postulate: it is the square root already present in the product. For a four-momentum \(p\) the grade-1 element \(\gamma(p)=\sum_{\mu}p_{\mu}\gamma^{\mu}\) satisfies \(\gamma(p)^2=Q(p)\), with \(Q(p)=-p_0^2+p_1^2+p_2^2+p_3^2\), both in the algebra and as Weyl matrices on \(\mathbb{C}^4\). The mass shell \(Q(p)=-m^2\) is the vanishing of \((\gamma(p)-im)(\gamma(p)+im)\). A spinor with \(\gamma(p)\Psi=im\Psi\) lies on that shell. In the chiral representation the slash is off-diagonal, so vanishing mass decouples the two Weyl blocks while a nonzero mass mixes them. Chirality \(\gamma^5\) reverses the sign of \(m\). Rest-frame solutions exist: a rest-frame particle spinor is assembled from a dual spinor \(u\) as the Weyl pair \((u,\,iu)\). Dual rotation of \(u\) leaves that rest-frame equation intact, so the internal \(\mathrm{SU}(2)\) of the dual sector is a symmetry of the mass shell at \(p=(m,0,0,0)\). The two Weyl blocks are orthogonal in \(\mathbb{C}^4\), and the inner product of two such rest-frame solutions that differ by a dual rotation is twice the dual-rotor amplitude,
\[
\langle\Psi(u)\mid\Psi\bigl(R(\beta)u\bigr)\rangle_{\mathbb{C}^4}
  =2\,\langle u\mid R(\beta)u\rangle_{\mathbb{C}^2}.
\]
The spinor matrix element of the dual rotor is therefore not an extra object placed beside the Dirac equation. It is the Hilbert pairing of rest-frame solutions related by an internal dual rotation. That pairing is still not \(e^{iEt/\hbar}\), and it does not produce a spacetime PDE.

This is the free Dirac equation in momentum space, equivalently the plane-wave form of \((i\gamma^{\mu}\partial_{\mu}-m)\Psi=0\). Squaring the slash recovers the Klein--Gordon symbol \(Q(p)+m^2\). The first-order structure is inherited from the Clifford product, not from a linearisation of a dual-rotor PDE about a torsional vacuum. Identifying \(m\) with dual-rotor stiffness remains a geometric reading, not a derived equality. The time gamma \(\gamma^0\) is the Minkowski vector \(e_0\), not the hyperbolic bivector used as a spin projector.

\section{Wave equations and geometric readings}

A slowly varying dual-rotor background may be read as suggesting the familiar laboratory wave equations
\begin{align}
i\hbar\frac{\partial\psi}{\partial t}
  &=\Bigl[-\frac{\hbar^2}{2m}\nabla^2+V(\mathbf{x})\Bigr]\psi, \\
\square\psi+\frac{m^2c^2}{\hbar^2}\psi&=0.
\end{align}
The passage from the \(\mathrm{SU}(2)\) kinematics, or from the affine mismatch, to these scalar PDEs is a heuristic expansion, not a derived theorem. Non-relativistic and relativistic quantum mechanics remain available as low-energy languages of dual-rotor dynamics; they are not forced by the character or by the spinor matrix element.

Wave-function collapse, in the same register, is readable as forced alignment of a particle's dual rotor to a macroscopic observer whose many dual rotors are locked into a stiff collective state. Apparent randomness is then ignorance of the initial dual angles, six real numbers per particle. Because the dual sector is compact, the realignment is instantaneous in the dual dimensions. This is a geometric picture of the measurement postulate, not a derivation of it.

The same 16-dimensional network is readable as transmitting macroscopic gravity through averaged mismatch and microscopic interference through local dual-rotor geometry. The diagnostic \(J\) of this note is not the parameter-space Killing quadratic of the main paper. Gravity and quantum non-locality may be the same algebra at different scales; that remains a reading.

\section{Relation to de Broglie--Bohm mechanics}

The pilot-wave interpretation of de Broglie and Bohm posits an external guiding field that steers particles along deterministic trajectories~\cite{bohm1952}. In the dual-spacetime framework no external field is required. The dual rotor that appears in the observer overlap is the particle's own dual rotor, evaluated at a hypothetical location of its centre of mass. The hidden variables, if one wishes to speak that way, are the six real rotor angles. The ``pilot wave'' is the particle viewed from the dual sector.

Bohmian mechanics still needs a scalar wave function on spacetime and a quantum potential. The dual-rotor geometry supplies an \(\mathrm{SU}(2)\) kinematics and a relative-rotor pairing; it does not, by itself, reconstruct the quantum potential or prove a Bell violation~\cite{bell1964}. Compactness of the dual sector makes non-locality available as a geometric option, not as a theorem of the character.

\section{Experimental signatures}

The identifications that remain readings still yield distinguishable expectations.

First, matter-wave interferometry would see a weak dependence of fringe visibility on the collective torsional stiffness of nearby masses, and a phase shift linear in the angular velocity of a large rotating mass---a gravitational Faraday-like effect from dual-rotor frame dragging.

Second, high-frequency gravitational-wave detectors operating near Compton frequencies would see resonant features only if a harmonic dual-rotor model is physically realised. The working oscillator, if it is the right model, sits at \(\sqrt{2m}\), not at \(m\). Absence of a narrow line at either frequency would constrain the dynamics, not the dual-rotor kinematics.

Third, if collapse is dual-rotor synchronisation, repeated weak measurements within a rotor relaxation time would show small systematic deviations from the Born rule. Such experiments lie within the reach of superconducting qubits or trapped ions.

These signatures distinguish a dual-spacetime origin of matter waves from standard quantum mechanics and from hidden-variable theories that retain an external pilot field. They do not test the character identity or the spinor matrix element, which are algebraic.

\section{Conclusions}

The matter wave is not the trace of a dual rotor. That trace is real because it sums two opposite half-angle phases. What the dual sector actually supplies is an \(\mathrm{SU}(2)\) kinematics on \(\mathbb{C}^2\): a real character, a complex spinor matrix element whose modulus is the polar tilt of the dual axis against the observer, and a relative-rotor pairing that compares two dual rotors without leaving that group. Off-axis the pairing's character records axis alignment; the noncommutative remainder is traceless. The same matrix element is the overlap of rest-frame Dirac solutions that differ by an internal dual rotation. The scalar de Broglie phase is a further identification, available as linear lag in the written particle action, not as a harmonic motion derived from rest mass, and not as the first-order Dirac factor already present in \(\mathrm{Cl}(3,1)\).

Quantum mechanics, in this geometry, is not an extra Hilbert space placed beside gravity. It is the dual-rotor kinematics of the same 16-dimensional algebra. The kinematics is sharp. The dynamical bridge to laboratory wave equations is not.

\begin{thebibliography}{9}

\bibitem{dst-main}
Anonymous,
``Gravity as Torsion between Dual Spacetime: A Biquaternionic Reformulation of General Relativity'' (2025).
\url{https://github.com/hypernumbernet/dual-spacetime-doc/blob/main/double-spacetime-theory.pdf}.

\bibitem{dst-dynamic-equation}
Anonymous,
``Dual-Rotor Dynamics in Double Spacetime Theory: Written Action, Free Mismatch, and a Neighbouring Oscillator'' (2025).
\url{https://github.com/hypernumbernet/dual-spacetime-doc/blob/main/dst-dynamic-equation.pdf}.

\bibitem{debroglie1924}
L.~de Broglie,
``Recherches sur la th\'eorie des quanta,''
\textit{Ann.\ Phys.} (Paris) \textbf{3}, 22--128 (1925).

\bibitem{bohm1952}
D.~Bohm,
``A suggested interpretation of the quantum theory in terms of ``hidden'' variables.\ I,''
\textit{Phys.\ Rev.} \textbf{85}, 166--179 (1952).

\bibitem{bell1964}
J.~S.~Bell,
``On the Einstein Podolsky Rosen paradox,''
\textit{Physics Physique Fizika} \textbf{1}, 195--200 (1964).

\end{thebibliography}

\end{document}
